\section{Introduction}\label{sec:introduction}

Let \(n\geq8\) be even and write \(G_n=C_n^2=C_n(1,2)\) for the graph on
\(\mathbb Z/n\mathbb Z\) whose edges join vertices at cyclic distance one or
two. An edge signing is a map
\(\sigma:E(G_n)\to\{\pm1\}\), and \(A_\sigma\) is its signed adjacency
matrix. We study
\[
m_n=\min_{\sigma:E(G_n)\to\{\pm1\}}\rho(A_\sigma),
\]
focusing on a conjectured formula \(m_n=\rho_-(n)\) supplied by a
distinguished twisted signing. Our main result determines when this
formula holds. It does not determine \(m_n\) when the formula fails or
classify all minimizers when it holds. The distance-one and distance-two
edges are signed independently; this is not the graph obtained by squaring
a previously signed cycle. Switching conjugates \(A_\sigma\) by a diagonal
\(\{\pm1\}\)-matrix, so the problem naturally lives on switching classes
\cite{Zaslavsky1982}.

The question belongs to the fixed-graph signing problem highlighted by
Belardo, Cioab\u{a}, Koolen, and Wang
\cite{BelardoCioabaKoolenWang2019}. Signings also encode the new eigenvalues
of two-lifts \cite{BiluLinial2006}, and interlacing families give powerful
one-sided bounds, which become two-sided for bipartite graphs by spectral
symmetry \cite{MarcusSpielmanSrivastava2015}. The cycle squares considered
here are nonbipartite, and the issue is whether one prescribed signing
attains the two-sided optimum. The starting point is Suvagiya's
preprint \cite[Conjecture~3]{Suvagiya2026Signed}. It introduces the relevant
flux coordinates, identifies the twisted classes, computes their spectra,
and verifies the conjectured equality through order \(18\). We retain those
inputs and determine the truth set.

The answer is nonmonotone. Equality holds at every even order through \(30\),
fails at \(32\), returns for \(34,36,38\), fails again at \(40\), and returns
for \(42,44,46\). Failure then persists at every even order from \(48\)
onward. The pattern is summarized in Table~\ref{tab:truth-pattern}; in
particular, the first counterexample does not predict the eventual regime.

\begin{table}[t]
\centering
\small
\begin{tabular}{c@{\qquad}cccccc}
\toprule
even order \(n\)&\(8\text{--}30\)&\(32\)&\(34\text{--}38\)&\(40\)&\(42\text{--}46\)&\(n\geq48\)\\
\midrule
comparison&\(m_n=\rho_-(n)\)&\(m_n<\rho_-(n)\)&\(m_n=\rho_-(n)\)&\(m_n<\rho_-(n)\)&\(m_n=\rho_-(n)\)&\(m_n<\rho_-(n)\)\\
\bottomrule
\end{tabular}
\caption{Truth pattern for the conjectured equality; interval columns contain
only even orders.}
\label{tab:truth-pattern}
\end{table}

Set
\[
\rho_-(n)^2=4+2\cos\frac{2\pi}{n}+2\cos\frac{4\pi}{n},
\qquad \theta_n=\rho_-(n)^2.
\]
An antiperiodic Fourier calculation shows that the twisted class attains
\(\rho_-(n)\) for every even \(n\geq8\), hence \(m_n\leq\rho_-(n)\). The
reverse inequality at a valid order is a separate global conclusion.

\begin{theorem}[Classification of the conjectured equality]
\label{thm:complete-classification}
For every even \(n\geq8\), the distinguished twisted signing attains
\(\rho_-(n)\), and
\[
m_n<\rho_-(n)
\quad\Longleftrightarrow\quad
n=32,\quad n=40,\quad\text{or}\quad n\geq48.
\]
Equivalently,
\[
m_n=\rho_-(n)
\quad\Longleftrightarrow\quad
n\in\{8,10,12,14,16,18,20,22,24,26,28,30,34,36,38,42,44,46\}.
\]
\end{theorem}

Two proof components have different roles. Finite classification resolves
the irregular small-order pattern. A structural argument explains why
failure is eventually uninterrupted: a low-spectral periodic phase supports
localized interfaces, and well-separated copies close on all sufficiently
large rings. The conjunction of the analytic tail and the finite bridge gives
the sharp onset at \(48\).

The comparison phase has period eight and is distinct from the twisted
attaining class. In its quadrilateral-flux word, positive sites occur at
spacing four. Replacing one spacing by a gap \(g\) creates a phase slip with
excess charge \(q=g-4\). On a ring, the charges satisfy a mod-eight closing
law, whereas each charge also transports a local translation sector modulo
four. These are separate constraints. The six-gap phase slip \(G_6\) has
charge two, so one, two, or three separated copies close the three nonzero
even residue classes modulo eight.

The reference phase has upper squared spectral edge
\[
\eta=4+\sqrt{10+2\sqrt5}<8.
\]
The central local theorem identifies \(G_6\) as the elementary low-cost
interface above this bulk. For the bilateral positive single-gap operator
\(A_g\), put \(H_g=A_g^2\) and
\(E(g)=\sup\operatorname{Spec}(H_g)\).

\begin{theorem}[Reference and single-gap spectral hierarchy]
\label{thm:single-gap-hierarchy}
There is an exactly certified level \(c_6\in(\eta,8)\) such that, for both
lifts and orientations,
\[
E(4)=\eta<c_6=E(6),
\qquad
\operatorname{Spec}_{\mathrm{ess}}(H_6)
=\operatorname{Spec}(H_{\mathrm{ref}}),
\qquad
\dim\ker(H_6-c_6)=2.
\]
Moreover, for every positive integer \(g\notin\{4,6\}\),
\[
E(g)>c_6+\frac1{250}.
\]
\end{theorem}

Thus \(G_6\) uniquely minimizes the edge among abnormal positive single
gaps; no claim is made about arbitrary multi-gap or finite-core interfaces.
The defining polynomial, physical matching argument, and rational isolation
of \(c_6\) appear in Section~\ref{sec:elementary-g6} and
Appendix~\ref{app:g6-certification}. Patch identification transfers the
reference and \(G_6\) bounds to finite rings. A discrete IMS partition then
controls the squared spectral top between patches and proves failure for all
even \(n\geq240\). Finite certificates cover \(48\leq n<240\), while
switching-class and finite-state arguments settle the smaller orders.

Several established strands meet in this proof. Signed-cycle spectra and
fixed-order signed extremal problems provide the closest spectral context
\cite{AkbariBelardoDodongehNematollahi2018,
BelardoBrunettiCiampella2021,BrunettiStanic2022}. Circulant and
skew-circulant Fourier analysis is classical
\cite{Davis1979,KarnerSchneidUeberhuber2003}, and Floquet methods on periodic
discrete graphs underlie the reference and interface calculations
\cite{Kuchment1993,KorotyaevSaburova2017,KorotyaevSaburova2023,
KuchmentVainberg2006}. The localization step is a finite-range discrete
instance of the IMS principle \cite{CyconFroeseKirschSimon1987}. For the
finite part, necklace and bracelet generation
\cite{FredricksenMaiorana1978,Sawada2001} and proof-producing exhaustive
classification \cite{Lam1991,GoedgebeurSchaudt2018,Goedgebeur2020} provide
methodological precedents. Combining an analytic large-order argument with
finite closure also occurs in all-order graph classifications
\cite{LinNing2021,GoedgebeurRendersWienerZamfirescu2024}. The new content is
the truth set for the
twisted candidate and the \(G_6\) mechanism behind eventual failure. To the
best of our knowledge, no previous work determines, for every even \(n\),
whether a prescribed signing minimizes spectral radius among all signings of
\(C_n(1,2)\).

Unless stated otherwise, every finite computation used in a proof is carried
out in integer, rational, or certified algebraic arithmetic under the
protocol of Section~\ref{sec:finite-completion} and the supplement. The
supplement identifies which objects are reconstructed independently and
which symbolic routines are shared. Section~\ref{sec:switching-reference}
sets up switching coordinates, the twisted candidate, and the period-eight
reference. Section~\ref{sec:gaps-charges-sectors} develops gaps, charges,
and sector transport. Sections~\ref{sec:elementary-g6} and
\ref{sec:single-gap} prove the local hierarchy. Section~\ref{sec:finite-rings}
constructs and localizes finite rings, and
Section~\ref{sec:finite-completion} closes the remaining orders. The
appendices give the human-checkable algebraic certificate and the
finite-classification completeness arguments; artifact-level data and
verification commands are in the supplement.
